\section{\label{III_1_C} Explicabilité et transparence}  

Le contrôle de la structure et de la qualité des données dépend également de la responsabilité de l'éditeur logiciel via-à-vis d'une transparence et d'une explicabilité du processus mis en place. Comme le rapelle Jonas Arnold, la gestion de la qualité de l'indexation dépend des utilisateurs, ainsi "la transparence envers ces derniers est primordiale." \footcite{zotero-item-584}. Afin de garantir cette transparence, les processus automatiques implémentés pourront être expliqués à deux niveaux au sein de l'interface. 

Tout d'abord, via une note d'information en amont du déclenchement du processus, au sein du diagramme de décision schématisé précédemment (\ref{fig:validation}). Avant le déclenchement du processus, l'interface devra indiquer que la fonctionnalité utilisée est fondée sur l'intelligence artificielle. Le même principe est prévu pour le développement de la recherche en langage naturel, via une infobulle qui en explique brièvement le fonctionnement à des fins de transparence (Voir Annexe B, \cref{figinfobulle}). De plus, une documentation pourra être mise à la disposition des utilisateurs sous la forme d'un mode d'emploi explicatif de la fonctionnalité. Skinsoft prévoit déjà ce type de documentation à l'usage des utilisateurs métier, via un manuel \textit{wiki} dédié à chaque espace de travail ets es fonctionnalités. Le mode d'emploi pourra mentionner de manière synthétique les critères d'extraction du modèle utilisé, justifier le choix des outils ainsi que la mention de leurs limites, pouvant guider les utilisateurs dans leur contrôle qualité des données. Il serait pertinent de mentionner également les limites liées à l'entraînement du modèle, qui apprend à extraire des informations brutes, mais qu'il n'apprend pas à les placer au sein d'une structure de données, en lien avec des normes métier. 

Il s'agira également de mentionner le périmètre au sein duquel sont analysées les données, qui est celui de l'éditeur seul. En effet, l'une des réticences majeures identifiés face à l'implémentation de fonctionnalités automatiques réside dans la perte de souveraineté des données patrimoniales lors de leur injection au sein de solutions commerciales externes \footcite{zotero-item-584}. Pour un éditeur comme Skinsoft, l'enjeu consiste à garantir une architecture logicielle éthique, étanche et souveraine, au sein de laquelle les données du client ne sont jamais réutilisées à son insu pour enrichir des réservoirs tiers.